\documentclass[12pt,a4paper]{article}

% Кодировка и языки
\usepackage[T1]{fontenc}
\usepackage[utf8]{inputenc}
\usepackage[russian]{babel}

% Геометрия страницы
\usepackage[left=2.5cm, right=2.5cm, top=2.5cm, bottom=2.5cm]{geometry}

% Для вставки кода/команд
\usepackage{listings}
\lstset{
  basicstyle=\ttfamily\small,
  breaklines=true,
  frame=single,
  rulecolor=\color{gray},
  showstringspaces=false,
  language=bash,
  aboveskip=10pt,
  belowskip=10pt
}

% Цвета для листингов
\usepackage{xcolor}

% Гиперссылки
\usepackage{hyperref}
\hypersetup{
  colorlinks=true,
  linkcolor=blue,
  urlcolor=blue
}

% Заголовок
\title{Система управления расписанием\\[0.3em]
\large Сетевая версия (TCP-сервер и клиент)}
\author{Бобнев Д.В.}
\date{\today}

\begin{document}
\maketitle
\tableofcontents
\newpage

% ----------------------------------------------------------------
\section{Руководство администратора}
\label{sec:admin}

\subsection{Системные требования}
Для сборки и запуска сервера и клиента необходимы:
\begin{itemize}
  \item Операционная система Linux / macOS (Windows поддерживается через WSL).
  \item Компилятор \texttt{g++} с поддержкой стандарта C++14 (или выше).
  \item Библиотека POSIX threads (\texttt{pthread}) – обычно входит в стандартную установку.
  \item Утилита \texttt{make} (опционально, для упрощения сборки).
\end{itemize}

\subsection{Состав дистрибутива}
Исходный код включает следующие файлы:
\begin{itemize}
  \item \texttt{shedule.h} – заголовочный файл с основными структурами и классами.
  \item \texttt{shedule.cpp} – реализация логики расписания (Database, парсер, match и др.).
  \item \texttt{DataGeneration.cpp} – генератор тестовых данных.
  \item \texttt{result.h}, \texttt{result.cpp} – класс Result, сериализация ответа.
  \item \texttt{server.cpp} – серверная часть (приём TCP‑соединений, обработка команд).
  \item \texttt{client\_lib.h}, \texttt{client\_lib.cpp} – библиотека клиента для подключения и отправки команд.
  \item \texttt{client.cpp} – консольный клиент.
\end{itemize}

\subsection{Сборка и установка}

\subsubsection{Сборка сервера}
Откройте терминал в каталоге с исходниками и выполните:
\begin{lstlisting}[language=bash]
g++ -std=c++14 -pthread -Wall -Wextra -o server \
  shedule.cpp result.cpp server.cpp DataGeneration.cpp
\end{lstlisting}

\subsubsection{Сборка клиента}
\begin{lstlisting}[language=bash]
g++ -std=c++14 -Wall -Wextra -o client \
  shedule.cpp result.cpp client_lib.cpp client.cpp
\end{lstlisting}

При необходимости можно добавить ключ \texttt{-O2} для оптимизации.

\subsection{Запуск и остановка}

\paragraph{Запуск сервера}
Сервер запускается из того же каталога, где находится файл данных \texttt{Data.txt} (или будет создан при первой генерации).
\begin{lstlisting}[language=bash]
./server
\end{lstlisting}
После запуска сервер начинает прослушивать TCP-порт 1234 и выводит сообщение о готовности.

\paragraph{Запуск клиента}
Клиент можно запускать на том же компьютере или удалённо, указав IP-адрес и порт (по умолчанию \texttt{127.0.0.1 1234}):
\begin{lstlisting}[language=bash]
./client [IP-адрес] [порт]
\end{lstlisting}

\paragraph{Остановка сервера}
Рекомендуемый способ – нажать \texttt{Ctrl+C} в терминале, где работает сервер. Это вызовет корректное завершение с освобождением всех ресурсов.  
Альтернативно можно отправить команду \texttt{SHUTDOWN} из клиента (если реализована), либо завершить процесс командой \texttt{killall server}.

\paragraph{Генерация начальных данных}
При первом запуске сервер пытается загрузить файл \texttt{Data.txt}. Если его нет, база будет пустой. Наполнить её можно командой \texttt{GENERATE} из клиента (см. \ref{sec:user}).

\subsection{Рекомендации по развёртыванию}
\begin{itemize}
  \item Сервер и клиент должны иметь доступ к одному и тому же каталогу с \texttt{Data.txt} (если используется относительный путь).
  \item Для многопользовательской работы сервер использует потоки; убедитесь, что в системе разрешено создание достаточного количества потоков.
  \item При необходимости изменить порт – отредактируйте константу в \texttt{server.cpp} и \texttt{client.cpp}.
\end{itemize}

% ----------------------------------------------------------------
\section{Руководство системного администратора}
\label{sec:sysadmin}

\subsection{Сетевой протокол}
Обмен данными между клиентом и сервером осуществляется поверх TCP. Порт по умолчанию – 1234. После установки соединения клиент отправляет запрос, сервер обрабатывает его и возвращает ответ. Соединение может оставаться открытым для нескольких последовательных команд (сеанс).

\subsection{Формат запроса (клиент $\rightarrow$ сервер)}
Запрос состоит из двух частей:
\begin{enumerate}
  \item \textbf{Заголовок}: 4 байта (uint32\_t) – общая длина тела запроса, переданная в сетевом порядке байт (big-endian, функция \texttt{htonl}).
  \item \textbf{Тело запроса}:
  \begin{itemize}
    \item 4 байта (uint32\_t) – идентификатор пользователя (UserID). Сервер может использовать это значение для сессий, но в актуальной реализации он заменяется на дескриптор сокета для уникальности.
    \item 4 байта (uint32\_t) – длина текстовой команды (без завершающего нуля).
    \item Последовательность байтов – текст команды в кодировке ASCII/UTF-8.
  \end{itemize}
\end{enumerate}

Пример (псевдокод):
\begin{lstlisting}[language=Python, caption=Пример формирования запроса]
cmd = "PRINT TEACHER == T_2"
uid = 1
body = pack("!I", uid) + pack("!I", len(cmd)) + cmd.encode()
header = pack("!I", len(body))
send(header + body)
\end{lstlisting}

\subsection{Формат ответа (сервер $\rightarrow$ клиент)}
Ответ также предваряется 4 байтами длины тела в big-endian. Тело ответа содержит сериализованный объект \texttt{Result} со следующими полями (порядок фиксирован):

\begin{enumerate}
  \item 1 байт – \texttt{success} (1 – успех, 0 – ошибка).
  \item 4 байта (uint32\_t) – длина строки \texttt{message}, затем сама строка.
  \item 4 байта (uint32\_t) – \texttt{affectedCount} (количество затронутых записей, например при REMOVE).
  \item 4 байта (uint32\_t) – количество записей \texttt{recordsCount}, затем каждая запись:
  \begin{itemize}
    \item Строка \texttt{subject} (4 байта длина + данные).
    \item Строка \texttt{teacher} (аналогично).
    \item 4 байта (uint32\_t) – \texttt{group}.
    \item 4 байта (uint32\_t) – \texttt{timeIndex}.
    \item 4 байта (uint32\_t) – \texttt{auditoryIndex}.
  \end{itemize}
  \item 4 байта – \texttt{rows}, 4 байта – \texttt{cols} (размеры матрицы).
  \item 4 байта – количество временных меток, затем каждая временная метка: строка в формате \texttt{"Weekday lesson\_num"} (например, \texttt{"Monday 1"}).
  \item 4 байта – количество аудиторий, затем строки с названиями аудиторий.
\end{enumerate}

Все целочисленные поля кодируются в big-endian (\texttt{htonl} при отправке, \texttt{ntohl} при чтении).

Строки передаются как последовательность: 4 байта длины + байты строки без завершающего нуля.

\subsection{Поддерживаемые команды}
Сервер поддерживает тот же набор команд, что и оригинальная консольная программа:
\begin{itemize}
  \item \texttt{SELECT <условия>} – выполнить поиск и сохранить выборку в сессию.
  \item \texttt{RESELECT <условия>} – уточнить текущую выборку.
  \item \texttt{PRINT} – показать полную матрицу расписания.
  \item \texttt{PRINT SELECT} – показать текущую выборку в виде матрицы.
  \item \texttt{PRINT TEACHER == <имя>} или \texttt{PRINT SUBJECT == <название>} – быстрый поиск с выводом матрицы.
  \item \texttt{ADD <день> <пара> <аудитория> <предмет> <преподаватель> <группа>}
  \item \texttt{REMOVE <условия>}
  \item \texttt{SAVE <файл> / LOAD <файл>}
  \item \texttt{GENERATE [параметры]} – генерация случайного расписания.
  \item \texttt{STATUS}, \texttt{CLEAR}, \texttt{HELP}, \texttt{EXIT}.
\end{itemize}
Полный синтаксис условий описан в пользовательском руководстве (\ref{sec:user}).

\subsection{Написание альтернативного клиента}
Для создания клиента на любом языке необходимо:
\begin{itemize}
  \item Установить TCP‑соединение с сервером.
  \item Реализовать описанную выше сериализацию запроса.
  \item Читать ответ (сначала 4 байта длины, затем тело) и десериализовать \texttt{Result}.
\end{itemize}
Учёт порядка байт (big-endian) обязателен. Игнорирование этого требования приведёт к неверной интерпретации чисел на платформах с little-endian.

% ----------------------------------------------------------------
\section{Руководство пользователя}
\label{sec:user}

\subsection{Запуск клиента}
После компиляции запустите клиент, указав при необходимости адрес и порт сервера:
\begin{lstlisting}[language=bash]
./client [server_ip] [port]
\end{lstlisting}
Если параметры опущены, используется \texttt{127.0.0.1} и порт \texttt{1234}. При успешном подключении появится приглашение \texttt{>}.

\subsection{Общие правила ввода команд}
\begin{itemize}
  \item Команды регистронезависимы (можно писать \texttt{print}, \texttt{SELECT} и т.д.).
  \item Параметры разделяются пробелами.
  \item Дни недели указываются на английском с заглавной буквы (например, \texttt{Monday}, \texttt{Tuesday}). Сокращения не допускаются.
  \item Числовые значения: номера уроков (1–6), группы, индексы.
  \item Строковые значения можно заключать в кавычки, но не обязательно, если нет пробелов.
\end{itemize}

\subsection{Основные команды}

\subsubsection{Просмотр расписания}
\begin{itemize}
  \item \texttt{PRINT} – выводит полную матрицу расписания. Строки – дни и пары, столбцы – аудитории. Ячейки содержат информацию о предмете, преподавателе и группе.
  \item \texttt{PRINT SELECT} – выводит текущую выборку (результат последнего \texttt{SELECT} или \texttt{RESELECT}).
  \item \texttt{PRINT TEACHER == <имя>} – выводит только те ячейки, где преподаватель совпадает с указанным. Например, \texttt{PRINT TEACHER == T\_2}.
  \item \texttt{PRINT SUBJECT == <название>} – аналогично для предмета.
\end{itemize}

\subsubsection{Поиск и фильтрация}
\begin{itemize}
  \item \texttt{SELECT <условия>} – выполняет поиск записей по заданным условиям и сохраняет их в сессию. Примеры:
  \begin{lstlisting}
SELECT GROUP == 5
SELECT TEACHER == T_3 AND DAY == Monday
SELECT SUBJECT == S_10 AND LESSON_NUM >= 3
  \end{lstlisting}
  Операторы: \texttt{==}, \texttt{!=}, \texttt{<}, \texttt{>}, \texttt{<=}, \texttt{>=}, \texttt{CONTAINS} (для строковых полей). Поля: \texttt{GROUP}, \texttt{TEACHER}, \texttt{SUBJECT}, \texttt{AUDITORY}, \texttt{DAY}, \texttt{DATE}, \texttt{TIME}, \texttt{LESSON\_NUM}. Условия объединяются словом \texttt{AND}.
  \item \texttt{RESELECT <условия>} – уточняет текущую выборку (применяет дополнительные условия). Если выборка пуста, ничего не делает.
\end{itemize}

\subsubsection{Добавление записи}
\begin{lstlisting}
ADD <день> <номер_пары> <аудитория> <предмет> <преподаватель> <группа>
\end{lstlisting}
Пример:
\begin{lstlisting}
ADD Tuesday 3 R_4 S_8 T_3 9
\end{lstlisting}
При конфликте (аудитория занята, преподаватель или группа уже имеют пару в это время) будет выведена ошибка «коллизия».

\subsubsection{Удаление записей}
\begin{lstlisting}
REMOVE <условия>
\end{lstlisting}
Условия задаются так же, как в \texttt{SELECT}. Команда удаляет все записи, удовлетворяющие условиям. Пример:
\begin{lstlisting}
REMOVE DAY == Saturday AND LESSON_NUM == 3
\end{lstlisting}

\subsubsection{Работа с файлами}
\begin{itemize}
  \item \texttt{SAVE <имя\_файла>} – сохраняет текущее расписание в текстовый файл (формат такой же, как у \texttt{Data.txt}).
  \item \texttt{LOAD <имя\_файла>} – загружает расписание из файла, полностью заменяя текущее содержимое базы.
\end{itemize}

\subsubsection{Генерация тестовых данных}
\begin{lstlisting}
GENERATE [уроков_в_день аудиторий групп преподавателей предметов заполненность]
\end{lstlisting}
Все параметры опциональны. Значения по умолчанию: 4 урока, 10 аудиторий, 10 групп, 8 преподавателей, 12 предметов, заполненность 0.6.  
Пример для большой базы:
\begin{lstlisting}
GENERATE 6 50 500 1000 500 0.8
\end{lstlisting}
После генерации автоматически выполняется \texttt{LOAD Data.txt}.

\subsubsection{Служебные команды}
\begin{itemize}
  \item \texttt{STATUS} – показывает количество временных слотов и аудиторий.
  \item \texttt{CLEAR} – очищает текущую выборку (сессию).
  \item \texttt{HELP} – выводит подсказку по командам.
  \item \texttt{EXIT} – завершает сеанс клиента и закрывает соединение. Сервер при этом продолжает работу.
\end{itemize}

\subsection{Особенности и ограничения}
\begin{itemize}
  \item В командах \texttt{PRINT} с условием поддерживаются только поля \texttt{TEACHER} и \texttt{SUBJECT} и только оператор \texttt{==}. Остальные комбинации вызовут ошибку с пояснением.
  \item Для \texttt{SELECT}, \texttt{RESELECT} и \texttt{REMOVE} операторы работают стандартно.
  \item Результаты \texttt{PRINT} отображаются в виде матрицы. Если записей много, возможно, потребуется расширить окно терминала для корректного отображения.
  \item Сессии разных клиентов изолированы друг от друга (каждый клиент имеет уникальный идентификатор сессии).
\end{itemize}

\subsection{Пример рабочей сессии}
\begin{lstlisting}[language=bash]
> GENERATE 6 50 500 1000 500 0.8
Сгенерировано и загружено
> STATUS
Слотов времени: 36, аудиторий: 50
> SELECT TEACHER == T_1
Выборка сохранена (14 записей)
> PRINT SELECT
(выводится матрица с 14 заполненными ячейками)
> ADD Monday 1 R_1 S_100 T_2 55
Добавлено
> PRINT TEACHER == T_2
(матрица с парами T_2)
> REMOVE GROUP == 55
Удалено записей: 1
> EXIT
Сеанс завершён
\end{lstlisting}

\end{document}